% ============================================================
% EV-PRAC-01 · ISR-401 Ingeniería de Software
% Compilador: pdflatex (TeX Live >= 2022) | Salida: 1 página A4
% ============================================================
\documentclass[10pt,a4paper]{extarticle}

\usepackage[utf8]{inputenc}
\usepackage[T1]{fontenc}
\usepackage{lmodern}
\usepackage[spanish]{babel}
\usepackage[top=0.7cm,bottom=0.7cm,left=1.0cm,right=1.0cm]{geometry}
\usepackage{enumitem}
\usepackage{booktabs}
\usepackage{tabularx}
\usepackage{array}
\usepackage[table]{xcolor}
\usepackage{microtype}
\usepackage{titlesec}
\usepackage[hidelinks]{hyperref}

% ---------- Compactación ----------
\pagestyle{empty}
\setlength{\parindent}{0pt}
\setlength{\parskip}{1.0pt}
\setlength{\tabcolsep}{3.0pt}
\renewcommand{\arraystretch}{1.05}
\setlist{nosep,leftmargin=1.1em,itemsep=0pt,topsep=0pt,parsep=0pt}
\titleformat{\section}{\bfseries\small\color{blue!45!black}}{\thesection.}{0.4em}{}[\vspace{-4pt}\rule{\linewidth}{0.3pt}]
\titlespacing*{\section}{0pt}{2pt}{1pt}
\emergencystretch=2em

\newcommand{\req}[1]{\textit{``#1''}}
\newcommand{\bad}{\textcolor{red!70!black}{\textbf{Mala:}}}
\newcommand{\good}{\textcolor{green!45!black}{\textbf{Buena:}}}

\begin{document}

% ---------- Encabezado ----------
\begin{center}
  {\large\bfseries Fundamentos de la Ingenier\'ia de Requisitos}\\[1pt]
  {\footnotesize EV-PRAC-01 · ISR-401 Ingeniería de Software · \textbf{Kevin German Contreras Chavez} · Matrícula 1251167720 ·
  Repo: \url{https://github.com/kcontrerasc3-cloud/ISR401-KevinContreras/tree/main/EV-PRAC-01}}\\[1pt]
  {\scriptsize Bases: ISO/IEC/IEEE 29148:2018 · SWEBOK v4.0 (KA Software Requirements) · IREB CPRE Foundation Level v3.1}
\end{center}
\vspace{-4pt}

% ============================================================
\section{Definici\'on formal y propiedades}
Seg\'un \textbf{ISO/IEC/IEEE 29148:2018}, un \emph{requisito} es un enunciado que traduce o expresa una necesidad junto con sus restricciones y condiciones asociadas. \textbf{IREB} lo caracteriza como (i) una necesidad percibida por un stakeholder, (ii) una capacidad o propiedad que el sistema debe poseer, o (iii) la representaci\'on documentada de ambas. \textbf{SWEBOK v4.0} lo concibe como una propiedad que el software debe exhibir para resolver un problema del mundo real.

\textbf{Propiedades de un requisito individual (29148, cl.\,5.2.5):} necesario, apropiado, no ambiguo, completo, singular, factible, verificable, correcto y conforme (a plantilla y estilo). \textbf{Del conjunto de requisitos (cl.\,5.2.6):} completo, consistente, factible, comprensible y validable. Cada requisito debe llevar atributos de gesti\'on (identificador \'unico, fuente, prioridad, estado, riesgo) que habiliten su trazabilidad. Verbos normativos: \emph{deber\'a} (obligatorio), \emph{deber\'ia} (recomendado), \emph{podr\'a} (opcional).

% ============================================================
\section{Proceso de Ingenier\'ia de Requisitos}
Las actividades son iterativas e incrementales, no una secuencia r\'igida:
\begin{itemize}
  \item \textbf{Elicitaci\'on:} descubrir necesidades desde stakeholders, documentos y sistemas existentes mediante entrevistas, talleres, observaci\'on, cuestionarios, prototipos y escenarios; clasificaci\'on de necesidades (modelo de Kano).
  \item \textbf{An\'alisis:} detectar conflictos, omisiones y dependencias; modelado (casos de uso, UML, BPMN); evaluaci\'on de factibilidad, riesgos y negociaci\'on; priorizaci\'on (MoSCoW, AHP).
  \item \textbf{Especificaci\'on:} documentar en lenguaje natural con plantillas (p.\,ej. EARS: \emph{[Condici\'on] [Sujeto] [Acci\'on] [Objeto] [Restricci\'on]}) y/o modelos; produce StRS, SyRS o SRS y la l\'inea base.
  \item \textbf{Validaci\'on:} confirmar que los requisitos reflejan la necesidad real (validez) mediante revisiones, inspecciones, \textit{walkthroughs}, prototipos y derivaci\'on temprana de casos de prueba.
  \item \textbf{Gesti\'on:} control de cambios, versionado, matriz de trazabilidad (bidireccional), l\'ineas base y m\'etricas de volatilidad de requisitos.
\end{itemize}

% ============================================================
\section{Tipos de requisitos}
\footnotesize
\begin{tabularx}{\textwidth}{@{}>{\bfseries}p{2.2cm}XX@{}}
\toprule
Tipo & Definici\'on & Ejemplo \\ \midrule
Funcional & Servicio, comportamiento o transformaci\'on de datos que el sistema debe realizar. & \req{El sistema deber\'a permitir al usuario autenticado restablecer su contrase\~na mediante un enlace enviado por correo.} \\
No funcional (calidad) & Atributo de calidad o desempe\~no (ISO/IEC 25010: rendimiento, seguridad, usabilidad, fiabilidad, mantenibilidad). & \req{El 95\,\% de las b\'usquedas deber\'a responderse en $\le$\,2\,s con 500 usuarios concurrentes.} \\
De dominio & Derivado del \'ambito de aplicaci\'on: leyes, est\'andares y reglas de negocio; pueden ser funcionales o no. & \req{El sistema deber\'a calcular el IVA conforme a la normativa tributaria vigente.} \\
Restricci\'on & Limita las soluciones admisibles (tecnolog\'ia, costo, plazo, normativa, interfaz); no a\~nade funcionalidad. & \req{La base de datos deber\'a ser PostgreSQL 16 sobre Linux.} \\
\bottomrule
\end{tabularx}
\normalsize
\vspace{1pt}
{\scriptsize\textit{Nota:} IREB distingue tres clases (funcionales, de calidad y restricciones); los requisitos de dominio se clasifican dentro de ellas.}

% ============================================================
\section{Perspectiva de sistemas: StRS, SyRS y SRS}
\footnotesize
\begin{tabularx}{\textwidth}{@{}>{\bfseries}p{1.4cm}XXX@{}}
\toprule
 & StRS & SyRS & SRS \\ \midrule
Nivel & Necesidades de stakeholders (espacio del problema, perspectiva del negocio/usuario). & Sistema completo (hardware, software, personas, procesos) como caja negra. & Componente de software dentro del sistema (espacio de la soluci\'on). \\
Audiencia & Usuarios, clientes, patrocinadores. & Arquitectos de sistemas, ingenier\'ia, aseguramiento. & Desarrolladores, testers, mantenimiento. \\
Contenido & Misi\'on, operaci\'on, entorno, procesos y calidad esperada. & Funciones, desempe\~no, interfaces, modos de operaci\'on y restricciones del sistema. & Requisitos funcionales, de interfaz, calidad, datos y restricciones del software. \\
Relaci\'on & Fuente de la que se derivan los SyRS. & Deriva del StRS y se asigna a subsistemas. & Deriva del SyRS (software asignado) y base del dise\~no y de las pruebas. \\
\bottomrule
\end{tabularx}
\normalsize

% ============================================================
\section{Calidad de los requisitos}
\textbf{Criterios SMART:} \emph{Specific} (concreto, sin vaguedad), \emph{Measurable} (con m\'etrica y umbral), \emph{Achievable} (factible t\'ecnica y econ\'omicamente), \emph{Relevant} (trazable a una necesidad) y \emph{Time-bound} (con plazo o condici\'on temporal cuando aplica). \textbf{Verificabilidad:} m\'etodo finito y objetivo de comprobaci\'on (inspecci\'on, an\'alisis, demostraci\'on o prueba). \textbf{Consistencia:} ausencia de contradicciones y uso de glosario controlado. \textbf{No ambig\"uedad:} una sola interpretaci\'on posible; evitar vaguesas.

\smallskip
\footnotesize
\begin{tabularx}{\textwidth}{@{}>{\raggedright\arraybackslash}X>{\raggedright\arraybackslash}X@{}}
\toprule
\bad\ \textbf{Defecto} & \good \\ \midrule
\req{El sistema debe ser r\'apido y f\'acil de usar.} \newline \emph{No medible, ambiguo, no verificable.} &
\req{Con 200 usuarios concurrentes, el sistema deber\'a responder el 95\,\% de las b\'usquedas en $\le$\,2\,s.} \\ \midrule
\req{El sistema deber\'a soportar muchos usuarios y enviar reportes si es posible.} \newline \emph{No singular, vago, cl\'ausula de escape.} &
\req{El sistema deber\'a soportar 1\,000 sesiones concurrentes.} \req{El sistema deber\'a generar el reporte mensual en PDF en $\le$\,30\,s.} \\ \midrule
REQ-7: \req{La sesi\'on expirar\'a a los 15 min.} \newline REQ-12: \req{La sesi\'on nunca expirar\'a.} \emph{(Inconsistentes).} &
REQ-7: \req{La sesi\'on deber\'a cerrarse tras 15\,min de inactividad, reiniciando el temporizador con cada acci\'on.} \\
\bottomrule
\end{tabularx}
\normalsize

\vspace{2pt}
{\scriptsize\textbf{Referencias:} ISO/IEC/IEEE 29148:2018; IEEE SWEBOK Guide v4.0, 2024; IREB CPRE FL Handbook v3.1.}

\end{document}